% ===================================================================
%  ਸਾਰਣੀ ਨੰ. 7 — ਸਿਆਲ ਵਿਚ ਪਿੰਡ। — ਬਹਾਰ ਵਿਚ ਖੇਤ।
%  Traduction 2026 d'apres texto/io, controlee sur texto/fr, texto/en
%  et texto/hi. Consigne : voir texto/pa/10-sarni-01.tex.
% ===================================================================

%%K t07-noto-1 noto
\VUnotes{143.2mm}{%
(*) ਖਾਣੇ ਦੇ ਵੇਰਵੇ ਲਈ ਸਾਰਣੀਆਂ 6 ਤੇ 13 ਵੇਖੋ।%
}

%%K t07-apar-1 apar
\VUsaut{4.0mm}
\VUcentre{13.2pt}{90}{ਦੂਜੀ ਲੜੀ}
\VUsaut{3.0mm}
\VUfilet{16mm}
\VUsaut{5.6mm}
\VUtitre{15.0pt}{60}{ਸਾਰਣੀ ਨੰ. 7}
\VUsaut{3.0mm}

%%K t07-tit-1 sub
\VUcentre{12.0pt}{0}{\textit{ਪਹਿਲਾ ਦ੍ਰਿਸ਼।}}
\VUsaut{3.0mm}

%%K t07-tit-2 sub
\VUpk{10.2pt}{40}{ਸਿਆਲ ਵਿਚ ਪਿੰਡ।}
\VUsaut{4.0mm}

%%K t07-c1-01-1 p
1. --- ਨਵੇਂ ਵਰ੍ਹੇ ਦੀਆਂ ਛੁੱਟੀਆਂ ਵਿਚ ਮੈਂ ਆਪਣੇ ਪਿਤਾ ਨਾਲ ਨਵਾਂ-ਪਿੰਡ ਗਿਆ, ਇਕ ਨਿੱਕਾ
ਪਿੰਡ ਜਿੱਥੇ ਉਨ੍ਹਾਂ ਨੇ ਕੁਝ ਕੰਮ ਨਿਬੇੜਨਾ ਸੀ।

%%K t07-c1-02-1 p
2. --- ਅਸੀਂ ਉਹ \VUgras{ਡਾਕ ਗੱਡੀ} \textsuperscript{(11)} ਲਈ ਜੋ ਸ਼ਹਿਰ ਤੇ ਉਸ ਪਿੰਡ ਦੇ
ਵਿਚਕਾਰ ਚਲਦੀ ਹੈ; ਪਰ ਪਾਲਾ ਬਹੁਤ ਸੀ, ਤੇ ਇੰਨੀ ਔਖੀ ਗੱਡੀ ਵਿਚ ਸਫ਼ਰ ਰਤਾ ਵੀ ਸੋਹਣਾ ਨਾ
ਰਿਹਾ।

%%K t07-c1-03-1 p
3. --- ਇਸੇ ਲਈ ਸਾਨੂੰ ਸੱਚਮੁੱਚ ਖ਼ੁਸ਼ੀ ਹੋਈ ਜਦੋਂ ਅਸੀਂ \VUgras{ਗੱਡੀ ਵਾਲੇ} ਨੂੰ ਚੌਕ ਵਿਚ
\textit{ਚਿੱਟੇ ਰਿੱਛ} ਦੀ \VUgras{ਸਰਾਂ} \textsuperscript{(2)} ਦੇ ਸਾਹਮਣੇ ਗੱਡੀ ਰੋਕਦੇ
ਵੇਖਿਆ। \VUgras{ਸਰਾਂ ਵਾਲਾ} \textsuperscript{(4)} ਆਪਣੀ ਦਹਿਲੀਜ਼ ਉੱਤੇ ਸਾਡੀ ਉਡੀਕ ਕਰ
ਰਿਹਾ ਸੀ, ਚੁੱਪ-ਚਾਪ ਆਪਣੀ \VUgras{ਚਿਲਮ} ਪੀਂਦਾ ਹੋਇਆ।

%%K t07-c1-03-2 p
ਉਸ ਨੇ ਸਾਨੂੰ ਅੰਦਰ ਸੱਦਿਆ, ਤੇ ਚੁੱਲ੍ਹੇ ਵਿਚ ਚਟਕਦੀ ਵੇਲ ਦੀ ਅੱਗ ਦੇ ਸਾਹਮਣੇ ਬਹਿਣ ਲਈ
ਮੈਨੂੰ ਦੂਜੀ ਵਾਰ ਕਹਾਉਣ ਦੀ ਲੋੜ ਨਾ ਪਈ। ਫਿਰ ਮੈਨੂੰ ਉਸ ਕੰਬਾ ਦੇਣ ਵਾਲੀ \VUgras{ਉੱਤਰੀ
ਹਵਾ} ਦੀ ਕੋਈ ਫ਼ਿਕਰ ਨਾ ਰਹੀ ਜੋ ਪੁਰਾਣੇ \VUgras{ਬੋਰਡ} \textsuperscript{(3)} ਨੂੰ
ਚਰਚਰਾਉਂਦੀ ਸੀ ਤੇ ਚਿਮਨੀ ਤੋਂ ਨਿਕਲਦੇ \VUgras{ਧੂੰਏਂ} \textsuperscript{(1)} ਨੂੰ ਕਾਲੇ
ਵਰੋਲਿਆਂ ਵਿਚ ਉਡਾ ਲੈ ਜਾਂਦੀ ਸੀ।

%%K t07-c1-04-1 p
4. --- ਦੁਪਹਿਰ ਦੇ ਖਾਣੇ ਦਾ ਵੇਲਾ ਸੀ। ਤੇ ਹੋਰ ਦੇਰ ਕੀਤੇ ਬਿਨਾਂ ਅਸੀਂ ਉਸ ਸਾਦੇ ਪਰ ਸੁਆਦੀ
ਖਾਣੇ ਨਾਲ ਪੂਰਾ ਇਨਸਾਫ਼ ਕੀਤਾ ਜੋ ਸਾਨੂੰ ਵਰਤਾਇਆ ਗਿਆ (*)।

%%K t07-tit-3 sub
\VUpk{10.2pt}{40}{ਕੁਝ ਮੁਲਾਕਾਤਾਂ।}
\VUsaut{3.0mm}

%%K t07-c1-05-1 p
5. --- ਖਾਣੇ ਪਿੱਛੋਂ ਮੈਂ ਆਪਣੇ ਪਿਤਾ ਨਾਲ ਗਿਆ, ਜੋ ਬੀਮੇ ਦੇ ਏਜੰਟ ਹਨ, ਉਨ੍ਹਾਂ ਦੇ ਗਾਹਕਾਂ
ਨੂੰ ਮਿਲਣ। ਪਹਿਲਾਂ ਅਸੀਂ \VUgras{ਲੁਹਾਰ} \textsuperscript{(5)} ਕੋਲ ਗਏ, ਜੋ ਸਰਾਂ ਦੀ ਬਗ਼ਲ
ਵਿਚ ਰਹਿੰਦਾ ਹੈ। ਉਹ ਇਕ ਘੋੜੇ ਦਾ ਨਾਲ ਠੋਕਣ ਵਿਚ ਲੱਗਾ ਹੋਇਆ ਸੀ। ਮੈਂ ਉਸ ਦੇ ਸਹਾਇਕ ਦੇ ਜ਼ੋਰ
ਦੀ ਦਾਦ ਦਿੱਤੀ, ਜੋ ਘੋੜੇ ਦਾ \VUgras{ਸੁੰਮ} ਆਪਣੇ ਚਮੜੇ ਦੇ ਪਰਨੇ ਉੱਤੇ ਕੱਸ ਕੇ ਫੜੀ ਹੋਇਆ
ਸੀ, ਜਦੋਂ ਤਕ ਲੁਹਾਰ ਭਾਰੇ ਹਥੌੜੇ ਦੀਆਂ ਸੱਟਾਂ ਨਾਲ \VUgras{ਨਾਲ} \textsuperscript{(6)}
ਗੱਡਦਾ ਰਿਹਾ।

%%K t07-c1-06-1 p
6. --- ਫਿਰ ਅਸੀਂ ਪਿੰਡ ਦੇ \VUgras{ਮੋਚੀ} \textsuperscript{(7)} ਬੁੱਢੇ ਜੈਕਬ ਕੋਲ ਗਏ।
ਭਾਵੇਂ ਉਸ ਦੇ ਬੋਰਡ ਉੱਤੇ ਇਕ ਸ਼ਾਨਦਾਰ \VUgras{ਬੂਟ} ਬਣਿਆ ਹੋਇਆ ਸੀ, ਉਹ ਪੁਰਾਣੇ ਘਸੇ
ਜੁੱਤਿਆਂ ਵਿਚ ਨਵੇਂ ਤਲੇ ਲਾ ਕੇ ਕੰਮ ਚਲਾ ਰਿਹਾ ਸੀ।

%%K t07-c1-06-2 p
ਬੁੱਢੇ ਨੇ ਹੱਸਦੇ ਹੋਏ ਆਖਿਆ : « ਸ਼ਹਿਰ ਵਿਚ ਮੈਂ “ਜੁੱਤੇ ਬਣਾਉਣ ਵਾਲਾ” ਸੀ ਤੇ ਮੇਰੇ ਗਾਹਕ
ਹੀ ਨਾ ਸਨ। ਇੱਥੇ ਮੈਂ “ਮੋਚੀ” ਹਾਂ ਤੇ ਕੰਮ ਦੀ ਕੋਈ ਘਾਟ ਨਹੀਂ। »

%%K t07-c1-07-1 p
7. --- ਉਸ ਨੇ ਸਾਨੂੰ ਆਪਣੇ ਗੁਆਂਢੀ \VUgras{ਹਜਾਮ} \textsuperscript{(8)} ਦੀ ਗੱਲ ਦੱਸੀ,
ਜਿਸ ਦੇ ਗਾਹਕ ਸਾਰੇ ਬੇਤਸੱਲੀ ਰਹਿੰਦੇ ਹਨ। ਉਸ ਦੇ \VUgras{ਉਸਤਰੇ} ਹਮੇਸ਼ਾ ਖੁੰਢੇ ਰਹਿੰਦੇ ਹਨ
ਤੇ ਉਸ ਦੀ \VUgras{ਕੈਂਚੀ} ਕਦੇ ਠੀਕ ਨਹੀਂ ਕੱਟਦੀ। ਪਰ ਕਿਉਂਕਿ ਉਹੋ ਪਿੰਡ ਦਾ ਇੱਕੋ ਹਜਾਮ ਹੈ,
ਲੋਕਾਂ ਕੋਲ ਉਸੇ ਤੋਂ ਹਜਾਮਤ ਕਰਾਉਣ ਤੇ ਵਾਲ ਕਟਾਉਣ ਤੋਂ ਬਿਨਾਂ ਚਾਰਾ ਨਹੀਂ। ਉਹ ਉੱਤੋਂ ਕੰਜੂਸ
ਤੇ ਰੁੱਖਾ ਵੀ ਹੈ।

%%K t07-c1-08-1 p
8. --- ਖ਼ੁਸ਼ਕਿਸਮਤੀ ਨਾਲ ਦੂਜੇ \VUgras{ਗੁਆਂਢੀ} ਉਸ ਵਰਗੇ ਨਹੀਂ।
\VUgras{ਗੱਡੀ ਬਣਾਉਣ ਵਾਲਾ} \textsuperscript{(9)}, ਮਿਸਾਲ ਲਈ, ਭਲਾ ਬੰਦਾ ਹੈ, ਮਿਹਨਤੀ ਤੇ
ਉਪਕਾਰੀ। ਉਸ ਨੂੰ ਗੱਡੀ ਮੁਰੰਮਤ ਨੂੰ ਦੇਣਾ ਸੁਖ ਦੀ ਗੱਲ ਹੈ। ਉਸ ਦਾ ਕੰਮ ਹਮੇਸ਼ਾ ਸਾਫ਼-ਸੁਥਰਾ
ਹੁੰਦਾ ਹੈ।

%%K t07-c1-09-1 p
9. --- ਬੁੱਢੇ ਜੈਕਬ ਨੂੰ \VUgras{ਕਾਠੀ-ਸਾਜ਼} \textsuperscript{(10)} ਤੋਂ ਵੀ ਕੋਈ ਸ਼ਿਕਾਇਤ
ਨਾ ਸੀ, ਜਿਸ ਦੀ ਉਹ ਕੰਮ ਦੇ ਦਬਾਅ ਵਿਚ ਕਦੇ ਕਦੇ \VUgras{ਸਾਜ਼} ਸੀਉਣ ਵਿਚ ਮਦਦ ਕਰਦਾ ਹੈ।

%%K t07-c1-09-2 p
ਮੇਰੇ ਪਿਤਾ ਨੇ ਉਸ ਕੋਲੋਂ ਉਸ ਦੇ ਪਰਿਵਾਰ ਬਾਰੇ ਪੁੱਛਿਆ। « ਸਾਰੇ ਰਾਜ਼ੀ-ਬਾਜ਼ੀ ਹਨ। ਮੇਰੀ
ਪੋਤਰੀ \VUgras{ਸਕੂਲ} \textsuperscript{(12)} ਵਿਚ ਹੈ। ਉਸ ਦੀ \VUgras{ਅਧਿਆਪਕਾ}
\textsuperscript{(13)} ਉਸ ਤੋਂ ਬਹੁਤ ਖ਼ੁਸ਼ ਹੈ; ਉਹ ਚੰਗੀ \VUgras{ਵਿਦਿਆਰਥਣ}
\textsuperscript{(14)} ਹੈ। ਉਹ ਮੇਰੇ ਉਸ ਸ਼ਰਾਰਤੀ ਪੁੱਤਰ ਵਰਗੀ ਨਹੀਂ; ਉਸ ਨੂੰ ਖੇਡ ਤੋਂ
ਬਿਨਾਂ ਕੁਝ ਨਹੀਂ ਸੁੱਝਦਾ। \VUgras{ਮੁੱਖ ਅਧਿਆਪਕ} \textsuperscript{(15)} ਉਸ ਨੂੰ ਕੁਝ ਵੀ
ਨਹੀਂ ਸਿਖਾ ਸਕਦੇ। »

%%K t07-tit-4 sub
\VUsaut{3.0mm}
\VUpk{10.2pt}{40}{ਪਿੰਡ ਦੀਆਂ ਗੱਲਾਂ।}
\VUsaut{3.0mm}

%%K t07-c1-10-1 p
10. --- ਫਿਰ ਬੁੱਢੇ ਨੇ ਸਾਨੂੰ ਪਿੰਡ ਦੀਆਂ ਖ਼ਬਰਾਂ ਸੁਣਾਈਆਂ। ਇਕ ਨਵਾਂ ਸਕੂਲ ਬਣਨਾ ਸੀ,
ਕਿਉਂਕਿ \VUgras{ਪਿੰਡ ਦੇ ਘਰ} ਵਾਲਾ ਬਹੁਤ ਨਿੱਕਾ ਪੈ ਗਿਆ ਸੀ। ਇਹ ਸੌਖਾ ਹੁੰਦਾ, ਕਿਉਂਕਿ
\VUgras{ਚੱਕੀ ਵਾਲੇ} ਦੇ ਬਾਗ਼ ਦਾ ਇਕ ਹਿੱਸਾ ਖ਼ਰੀਦ ਲੈਣਾ ਹੀ ਬਹੁਤ ਸੀ। ਵਿਚਾਰੇ ਲਈ ਇਹ ਬੜਾ
ਚੰਗਾ ਰਹਿੰਦਾ। ਜਦੋਂ ਦਾ ਪਾਲਾ ਪਿਆ, \VUgras{ਚੱਕੀ} \textsuperscript{(16)} ਦੇ
\VUgras{ਪੁੜ} ਕਣਕ ਨਾ ਪੀਹ ਸਕੇ, ਕਿਉਂਕਿ \VUgras{ਪਹੀਆ} \textsuperscript{(17)}
\VUgras{ਬਰਫ਼} \textsuperscript{(25)} ਵਿਚ ਫਸ ਗਿਆ ਸੀ, ਜੋ \VUgras{ਨਾਲੇ}
\textsuperscript{(23)} ਦੀ ਸੀ।

%%K t07-c1-11-1 p
11. --- « ਬਣਤਰ ਦੀ ਗੱਲ ਚੱਲੀ ਤਾਂ — ਸਾਡਾ ਨਵਾਂ \VUgras{ਗਿਰਜਾ} \textsuperscript{(18)}
ਵੇਖਿਆ ਜੇ? ਬਰਫ਼ ਦੀ ਚਾਦਰ ਦੇ ਹੇਠਾਂ ਉਹ ਬੜਾ ਸੋਹਣਾ ਲਗਦਾ ਹੈ। \VUgras{ਕਾਂ}
\textsuperscript{(19)}, ਜੋ ਆਪਣਾ ਪੁਰਾਣਾ ਘੰਟਾ-ਘਰ ਹੁਣ ਪਛਾਣਦੇ ਹੀ ਨਹੀਂ, ਉਦਾਸ ਹੋ ਕੇ
\VUgras{ਕਬਰਸਤਾਨ} \textsuperscript{(20)} ਦੇ ਵੱਡੇ ਰੁੱਖ ਉੱਤੇ ਬੈਠੇ ਰਹਿੰਦੇ ਹਨ। »

%%K t07-c1-12-1 p
12. --- ਕਬਰਸਤਾਨ ਦਾ ਨਾਂ ਆਉਂਦੇ ਹੀ ਮੋਚੀ ਨੂੰ ਉਹ ਲੋਕ ਯਾਦ ਆ ਗਏ ਜਿਨ੍ਹਾਂ ਨੂੰ ਮੇਰੇ ਪਿਤਾ
ਜਾਣਦੇ ਸਨ ਤੇ ਜੋ ਪਿਛਲੇ ਵਰ੍ਹੇ ਤੁਰ ਗਏ। « ਕਿੰਨੀਆਂ \VUgras{ਕਬਰਾਂ}
\textsuperscript{(21)}, ਮੇਰੇ ਭਲੇ ਜਨਾਬ, \VUgras{ਸਰੂ} \textsuperscript{(22)} ਦੇ ਹੇਠਾਂ
ਹੋਰ ਜੁੜ ਗਈਆਂ! ਮੌਤ ਸਾਡੇ ਬੁੱਢੇ ਮੁਨਸ਼ੀ ਨੂੰ ਲੈ ਗਈ। ਸਾਰਾ ਪਿੰਡ ਉਨ੍ਹਾਂ ਦੇ
\VUgras{ਸਸਕਾਰ} ਵਿਚ ਸੀ। »

%%K t07-c1-13-1 p
13. --- « ਉਹ ਰਹੇ ਉਨ੍ਹਾਂ ਦੇ ਵਾਰਸ, ਨਾਲੇ ਉੱਤੇ ਬਰਫ਼ ਤਿਲਕ ਰਹੇ ਹਨ। ਥੋੜੀ ਦੇਰ ਪਹਿਲਾਂ ਹੀ
ਆਏ ਸਨ ਕਿ ਮੈਂ ਉਨ੍ਹਾਂ ਦੀ \VUgras{ਬਰਫ਼-ਜੁੱਤੀ} \textsuperscript{(26)} ਦਾ ਤਸਮਾ ਜੋੜ
ਦੇਵਾਂ, ਜੋ ਟੁੱਟ ਗਿਆ ਸੀ। »

%%K t07-c1-14-1 p
14. --- « ਤੇ ਉਹ ਰਿਹਾ, ਨਾਲੇ ਦੇ ਕੰਢੇ, ਨਵਾਂ \VUgras{ਪਿੰਡ ਦਾ ਚੌਕੀਦਾਰ}
\textsuperscript{(27)}। ਉਹ ਪੁਰਾਣਾ ਸਿਪਾਹੀ ਹੈ ਤੇ ਪਿੰਡ ਦੇ ਸ਼ਰਾਰਤੀ ਮੁੰਡਿਆਂ ਦਾ ਕਾਲ। ਉਸ
ਦੇ \VUgras{ਪਟੇ} \textsuperscript{(29)} ਤੇ ਉਸ ਦੀ \VUgras{ਝਾਲਰ ਵਾਲੀ ਟੋਪੀ}
\textsuperscript{(28)} ਦਿਸਦੇ ਹੀ ਉਹ ਨੱਸ ਖੜਦੇ ਹਨ। »

%%K t07-c1-15-1 p
15. --- « ਓਹੋ! ਉਹ ਰਿਹਾ ਬੁੱਢਾ ਜੋਜ਼ਫ਼, ਨਾਨਬਾਈ ਲਈ \VUgras{ਲੱਕੜਾਂ ਦੀਆਂ ਗੱਠੜੀਆਂ ਦੀ
ਗੱਡੀ} \textsuperscript{(30)} ਲਿਆਉਂਦਾ ਹੋਇਆ। --- ਠੀਕ ਆਖਿਆ, ਮੇਰੇ ਪਿਤਾ ਬੋਲੇ, ਮੈਂ ਉਸ
ਦੇ \VUgras{ਖੱਚਰਾਂ} \textsuperscript{(31)} ਦੀ ਜੋੜੀ ਪਛਾਣਦਾ ਹਾਂ। ਮੈਨੂੰ ਉਸ ਨਾਲ ਇਕ ਗੱਲ
ਕਰਨੀ ਹੈ, ਇਹੋ ਮੌਕਾ ਹੈ। ਨਮਸਕਾਰ, ਬੁੱਢੇ ਜੈਕਬ। »

%%K t07-c1-16-1 p
16. --- ਅਸੀਂ ਬਾਹਰ ਨਿਕਲ ਹੀ ਰਹੇ ਸੀ ਕਿ ਪਿੰਡ ਦੇ ਬੱਚੇ ਸਕੂਲ ਤੋਂ ਛੁੱਟ ਕੇ
\VUgras{ਬਰਫ਼ ਦੀ ਤਿਲਕ-ਪੱਟੀ} \textsuperscript{(33)} ਉੱਤੇ ਟੁੱਟ ਪਏ, ਇਸ ਖ਼ਤਰੇ ਨਾਲ ਕਿ
\VUgras{ਡਿੱਗਣ} \textsuperscript{(34)} ਉੱਤੇ ਹੱਥ-ਪੈਰ ਤੁੜਵਾ ਬੈਠਣ। ਉਨ੍ਹਾਂ ਵਿਚੋਂ ਇਕ ਨੇ
ਗੱਡੀ ਵਾਲੇ ਦੇ ਘਰੋਂ ਆਪਣੀ \VUgras{ਸਲੈੱਜ} \textsuperscript{(32)} ਲਿਆਂਦੀ, ਉਸ ਉੱਤੇ ਆਪਣਾ
ਇਕ ਸਾਥੀ ਬਿਠਾਇਆ, ਤੇ ਦੌੜ ਪਿਆ।

%%K t07-c1-17-1 p
17. --- ਦੂਜੇ ਇਕ \VUgras{ਬਰਫ਼ ਦੇ ਢੇਰ} \textsuperscript{(37)} ਵੱਲ ਲਪਕੇ, ਜੋ
\VUgras{ਪੁਲ} \textsuperscript{(40)} ਕੋਲ ਸੀ, ਤੇ ਉਸ \VUgras{ਬਰਫ਼ ਦੇ ਬੰਦੇ}
\textsuperscript{(39)} ਉੱਤੇ ਮਾਰਨ ਲੱਗੇ ਜੋ ਉਨ੍ਹਾਂ ਕੱਲ੍ਹ ਬਣਾਇਆ ਸੀ, ਜਿਸ ਨੂੰ ਉਨ੍ਹਾਂ ਇਕ
ਟੋਪੀ, ਇਕ ਚਿਲਮ ਤੇ ਮੋਢੇ ਉੱਤੇ ਲਟਕਿਆ ਇਕ ਬਸਤਾ ਦਿੱਤਾ ਹੋਇਆ ਸੀ।

%%K t07-c1-18-1 p
18. --- ਉਨ੍ਹਾਂ ਨੂੰ ਬਰਫ਼ੀਲੀ ਹਵਾ ਤੇ ਪਾਲੇ ਦੀ ਪਰਵਾਹ ਨਾ ਸੀ। ਉਨ੍ਹਾਂ ਵਿਚੋਂ ਬਹੁਤਿਆਂ ਕੋਲ
\VUgras{ਗਲੂਬੰਦ} \textsuperscript{(35)} ਤਕ ਨਾ ਸੀ, ਤੇ \VUgras{ਬਰਫ਼ ਦੇ ਗੋਲਿਆਂ}
\textsuperscript{(38)} ਦੀ ਲੜਾਈ ਪਿੱਛੋਂ ਨਿੱਘੇ ਹੋਣ ਲਈ ਮੁੱਠ ਭਰ ਤਪਦੇ
\VUgras{ਚੈਸਟਨਟ} \textsuperscript{(42)} ਹੀ ਉਨ੍ਹਾਂ ਨੂੰ ਬਹੁਤ ਸਨ, ਜੋ ਉਹ ਬੁੱਢੀ ਜੈਕਲੀਨ
\VUgras{ਵੇਚਣ ਵਾਲੀ} ਕੋਲੋਂ ਖ਼ਰੀਦਦੇ ਸਨ, ਜੋ ਆਪਣੀ ਬਹੁਤ ਪਤਲੀ \VUgras{ਚਾਦਰ}
\textsuperscript{(43)} ਵਿਚ ਪਾਲੇ ਨਾਲ ਕੰਬਦੀ ਰਹਿੰਦੀ ਹੈ।

%%K t07-tit-5 sub
\VUsaut{3.0mm}
\VUcentre{12.0pt}{0}{\textit{ਦੂਜਾ ਦ੍ਰਿਸ਼।}}
\VUsaut{3.0mm}

%%K t07-tit-6 sub
\VUpk{10.2pt}{40}{ਬਹਾਰ ਵਿਚ ਖੇਤ ਦਾ ਘਰ।}
\VUsaut{4.0mm}

%%K t07-c2-01-1 p
1. --- ਸਿਆਲ ਦੇ ਸਾਰੇ ਸੁਖਾਂ ਦੇ ਬਾਵਜੂਦ \VUgras{ਬਹਾਰ} ਦੇ ਮੁੜ ਆਉਣ ਦਾ ਅਸੀਂ ਡੂੰਘੇ
ਆਨੰਦ ਨਾਲ ਸੁਆਗਤ ਕਰਦੇ ਹਾਂ।

%%K t07-c2-01-2 p
ਮਈ ਦੇ ਮਹੀਨੇ ਵਿਚ ਸ਼ਾਮ ਨੂੰ ਖੇਤ ਦਾ ਵਿਹੜਾ ਕਿਹੋ ਜਿਹੀ ਜੋਸ਼ ਭਰੀ ਤਸਵੀਰ ਬਣਾਉਂਦਾ ਹੈ!

%%K t07-c2-02-1 p
2. --- ਦਿਸਹੱਦੇ ਉੱਤੇ \VUgras{ਸੂਰਜ} \textsuperscript{(1)} ਹੌਲੀ ਹੌਲੀ ਡੁੱਬ ਰਿਹਾ ਹੈ, ਤੇ
\VUgras{ਪਿੰਡ} \textsuperscript{(3)} ਨੂੰ ਆਪਣੀਆਂ ਲਾਲ \VUgras{ਕਿਰਨਾਂ}
\textsuperscript{(2)} ਨਾਲ ਭਰ ਰਿਹਾ ਹੈ। ਮਿਹਨਤੀ \VUgras{ਹਾਲੀ} \textsuperscript{(6)}
ਆਪਣਾ \VUgras{ਖੇਤ} \textsuperscript{(4)} ਵਾਹੁਣ ਦੀ ਕਾਹਲੀ ਵਿਚ ਹੈ।

%%K t07-c2-02-2 p
ਆਪਣੇ \VUgras{ਹਲ} \textsuperscript{(5)} ਨਾਲ, ਜਿਸ ਨੂੰ ਇਕ ਤਕੜਾ \VUgras{ਘੋੜਾ}
\textsuperscript{(7)} ਖਿੱਚਦਾ ਹੈ, ਉਹ \VUgras{ਸਿਆੜ} \textsuperscript{(8)} ਖਿੱਚਦਾ ਹੈ,
ਜਿਸ ਵਿਚ \VUgras{ਬੀ ਬੀਜਣ ਵਾਲਾ} \textsuperscript{(9)} ਮੁੱਠ ਮੁੱਠ ਉਹ \VUgras{ਬੀ}
ਛੱਟੇਗਾ ਜਿਸ ਤੋਂ ਸੁਨਹਿਰੀ ਸਿੱਟੇ ਨਿਕਲਣਗੇ।

%%K t07-c2-03-1 p
3. --- ਆਪਣੀਆਂ ਦਾਤਰੀਆਂ ਨਾਲ \VUgras{ਘਾਹ ਕੱਟਣ ਵਾਲਿਆਂ} \textsuperscript{(11)} ਨੇ
ਮੈਦਾਨ ਦਾ ਘਾਹ ਕੱਟ ਲਿਆ ਹੈ, ਸੁਕਾਉਣ ਵਾਲਿਆਂ ਨੇ \VUgras{ਸੁੱਕਾ ਘਾਹ}
\textsuperscript{(10)} ਨੂੰ \VUgras{ਪੰਜਿਆਂ} \textsuperscript{(12)} ਤੇ ਖੁਰਪਿਆਂ ਨਾਲ
ਪਰਤ-ਪਰਤ ਕੇ ਸੁਕਾ ਲਿਆ ਹੈ, ਤੇ ਹੁਣ ਉਹ ਮੁੜ ਰਹੇ ਹਨ।

%%K t07-c2-03-2 p
ਕੁਝ ਬਲਦਾਂ ਨਾਲ ਖਿੱਚੀ ਭਾਰੀ ਗੱਡੀ ਦੇ ਅੱਗੇ ਅੱਗੇ ਤੁਰਦੇ ਹਨ, ਮੋਢੇ ਉੱਤੇ ਆਪਣੇ ਔਜ਼ਾਰ ਚੁੱਕੀ।
ਦੂਜੇ, ਵੱਧ ਆਲਸੀ, ਮਹਿਕਦੇ ਸੁੱਕੇ ਘਾਹ ਉੱਤੇ ਆਰਾਮ ਨਾਲ ਪਸਰ ਗਏ ਹਨ।

%%K t07-c2-04-1 p
4. --- ਤੁਰਦੇ ਤੁਰਦੇ ਉਹ \VUgras{ਮਾਲੀ} \textsuperscript{(16)} ਨੂੰ ਦੋਸਤੀ ਨਾਲ ਨਮਸਕਾਰ
ਕਰਦੇ ਹਨ, ਜੋ \VUgras{ਬਾਗ਼} \textsuperscript{(13)} ਵਿਚ ਲੱਗਾ ਹੈ, \VUgras{ਫਲਾਂ ਵਾਲੇ
ਰੁੱਖਾਂ} ਦੀ ਛਾਂਗ ਕਰਦਾ ਤੇ \VUgras{ਨਾਸ਼ਪਾਤੀ ਦੇ ਰੁੱਖਾਂ} \textsuperscript{(14)} ਵਿਚ
ਕਲਮ ਲਾਉਂਦਾ। \VUgras{ਆਲੂਬੁਖ਼ਾਰੇ ਦੇ ਰੁੱਖ} \textsuperscript{(15)} ਫੁੱਲਾਂ ਨਾਲ ਲੱਦੇ ਹਨ,
ਤੇ ਚੰਗੀ ਰੂੜੀ ਪਾਏ \VUgras{ਸਬਜ਼ੀ ਦੇ ਬਾਗ਼} \textsuperscript{(17)} ਵਿਚ ਸਬਜ਼ੀਆਂ, ਗੋਭੀ ਤੇ
ਸਲਾਦ ਪਹਿਲਾਂ ਤੋਂ ਚੰਗੀਆਂ ਹੋ ਰਹੀਆਂ ਹਨ।

%%K t07-c2-04-2 p
ਨੀਵੀਂ ਕੰਧ ਉੱਤੇ ਇਕ ਸੋਹਣਾ \VUgras{ਮੋਰ} \textsuperscript{(18)} ਆਪਣੇ ਸ਼ਾਨਦਾਰ ਖੰਭ
ਵਿਖਾਉਣ ਨੂੰ ਪੂਛ ਖਿਲਾਰਦਾ ਹੈ।

%%K t07-tit-7 sub
\VUpk{10.2pt}{40}{ਖੇਤ ਦਾ ਵਿਹੜਾ।}
\VUsaut{3.0mm}

%%K t07-c2-05-1 p
5. --- \VUgras{ਤਬੇਲੇ} \textsuperscript{(19)} ਦੇ ਬੂਹੇ ਖੁੱਲ੍ਹੇ ਹਨ, ਤੇ ਖੁਰਲੀ ਤੇ
\VUgras{ਵਿਛਾਉਣ ਦੀ ਪਰਾਲੀ} \textsuperscript{(23)} ਵਿਖਾਈ ਦਿੰਦੀ ਹੈ ਉਸ \VUgras{ਘੋੜੀ}
\textsuperscript{(21)} ਦੀ, ਜਿਸ ਨੂੰ \VUgras{ਸਾਈਸ} \textsuperscript{(20)} ਨੇ ਹੁਣੇ
ਖੋਲ੍ਹਿਆ ਹੈ। \VUgras{ਵਛੇਰਾ} \textsuperscript{(22)}, ਆਪਣੀਆਂ ਲੰਮੀਆਂ ਲੱਤਾਂ ਉੱਤੇ ਇੰਨਾ
ਮਜ਼ੇਦਾਰ, ਛਾਲਾਂ ਮਾਰਦਾ ਆਪਣੀ ਮਾਂ ਦੇ ਪਿੱਛੇ ਤੁਰਦਾ ਹੈ।

%%K t07-c2-06-1 p
6. --- \VUgras{ਖੂਹ} \textsuperscript{(29)} ਕੋਲ \VUgras{ਗਾਂਵਾਂ}
\textsuperscript{(25)} ਉਸ ਲੱਕੜ ਦੀ \VUgras{ਨਾਂਦ} \textsuperscript{(26)} ਤੋਂ ਪਾਣੀ
ਪੀਣ ਜਾਂਦੀਆਂ ਹਨ ਜੋ ਉਨ੍ਹਾਂ ਦਾ ਘਾਟ ਹੈ। \VUgras{ਵੱਛਾ} \textsuperscript{(24)} ਇਸੇ ਮੌਕੇ
ਉੱਤੇ ਦੁੱਧ ਪੀ ਲੈਂਦਾ ਹੈ, ਤੇ \VUgras{ਬਲਦ} \textsuperscript{(27)}, ਚਾਰੇ ਦੀ ਚੰਗੀ ਵਾਸ਼ਨਾ
ਪਾ ਕੇ, ਖ਼ੁਸ਼ੀ ਨਾਲ ਅੜਿੰਗਦਾ ਹੈ ਤੇ ਆਪਣੇ ਤਕੜੇ \VUgras{ਸਿੰਗਾਂ} \textsuperscript{(28)}
ਵਾਲਾ ਸਿਰ ਮਾਣ ਨਾਲ ਚੁੱਕਦਾ ਹੈ।

%%K t07-c2-07-1 p
7. --- ਸਾਰੇ ਕੰਮ ਵਿਚ ਲੱਗੇ ਹਨ। \VUgras{ਨੌਕਰਾਣੀ} \textsuperscript{(32)} ਖੂਹ ਦਾ
\VUgras{ਰੱਸਾ} \textsuperscript{(31)} ਫੜੀ ਹੋਈ ਹੈ। ਉਹ ਡੰਗਰਾਂ ਨੂੰ ਪਿਆਉਣ ਲਈ ਇਕ
\VUgras{ਬਾਲਟੀ} \textsuperscript{(30)} ਵਿਚ ਪਾਣੀ ਖਿੱਚ ਰਹੀ ਹੈ। \VUgras{ਕੋਠੇ}
\textsuperscript{(33)} ਕੋਲ ਇਕ ਬੰਦਾ ਖਿੱਲਰੀ ਪਰਾਲੀ ਇਕੱਠੀ ਕਰ ਰਿਹਾ ਹੈ, ਤੇ
\VUgras{ਖੇਤ ਦੇ ਘਰ} \textsuperscript{(34)} ਦੇ ਬੂਹੇ ਦੇ ਸਾਹਮਣੇ ਇਕ ਬੁੱਢੀ
\VUgras{ਕੱਤਣ ਵਾਲੀ} \textsuperscript{(35)}, ਆਪਣੇ ਚਰਖ਼ੇ ਤੇ ਆਪਣੀ \VUgras{ਤੱਕਲੇ} ਨਾਲ,
ਪੁਰਾਣੇ ਦਿਨਾਂ ਵਾਂਗ \VUgras{ਸਣ} ਜਾਂ \VUgras{ਪਟਸਣ} ਕੱਤ ਰਹੀ ਹੈ।

%%K t07-tit-8 sub
\VUsaut{3.0mm}
\VUpk{10.2pt}{40}{ਖੇਤ ਦਾ ਮਾਲਕ।}
\VUsaut{3.0mm}

%%K t07-c2-08-1 p
8. --- \VUgras{ਖੇਤ ਦਾ ਮਾਲਕ} \textsuperscript{(36)} ਇਹ ਸਾਰਾ ਕੰਮ ਚਲਾਉਂਦਾ ਹੈ ਤੇ ਆਪਣੇ
ਬੰਦਿਆਂ ਨੂੰ ਹੁਕਮ ਦਿੰਦਾ ਹੈ। ਉਹ ਉਨ੍ਹਾਂ ਨੂੰ ਆਖਦਾ ਹੈ ਕਿ \VUgras{ਸੁਹਾਗਾ}
\textsuperscript{(40)} ਤੇ \VUgras{ਕਹੀ} \textsuperscript{(39)} ਢਕ ਕੇ ਰੱਖ ਦੇਣ, ਜੋ ਉਸ
\VUgras{ਘਰ} \textsuperscript{(38)} ਕੋਲ ਛੱਡ ਦਿੱਤੇ ਗਏ ਹਨ ਜੋ \VUgras{ਕੁੱਤੇ}
\textsuperscript{(37)} ਦਾ ਹੈ।

%%K t07-c2-09-1 p
9. --- ਉਸ ਦੀ ਹਾਕ ਨਾਲ ਇਕ \VUgras{ਕਬੂਤਰ} \textsuperscript{(41)} ਚੌਂਕ ਉੱਠਦਾ ਹੈ, ਜੋ
ਪੂਰੀ ਚਾਲ ਨਾਲ \VUgras{ਕਬੂਤਰਖ਼ਾਨੇ} \textsuperscript{(42)} ਵੱਲ ਉੱਡ ਜਾਂਦਾ ਹੈ, ਤੇ
\VUgras{ਕੁਕੜੀਆਂ} \textsuperscript{(43)} ਵੀ, ਜੋ ਛੇਤੀ ਛੇਤੀ \VUgras{ਡੱਬੇ}
\textsuperscript{(44)} ਵਿਚ ਮੁੜ ਜਾਂਦੀਆਂ ਹਨ।

%%K t07-c2-10-1 p
10. --- \VUgras{ਭੇਡਾਂ ਦੇ ਵਾੜੇ} \textsuperscript{(48)} ਦੇ ਸਾਹਮਣੇ, ਇਹ ਰਿਹਾ ਬੁੱਢਾ
\VUgras{ਆਜੜੀ} \textsuperscript{(45)} ਆਪਣਾ \VUgras{ਇੱਜੜ} \textsuperscript{(49)}
ਮੋੜਦਾ ਹੋਇਆ। ਆਪਣੀ \VUgras{ਲਾਠੀ} \textsuperscript{(46)} ਫੜੀ, ਜੋ ਉਸ ਦਾ ਸਾਥ ਓਨਾ ਹੀ ਨਹੀਂ
ਛੱਡਦੀ ਜਿੰਨਾ ਉਸ ਦਾ ਫਿੱਕਾ \VUgras{ਕੁੜਤਾ} \textsuperscript{(47)}, ਉਹ ਵਾੜੇ ਵੱਲ
\VUgras{ਛੱਤਰੇ} \textsuperscript{(50)}, \VUgras{ਭੇਡਾਂ} \textsuperscript{(53)},
\VUgras{ਲੇਲੇ} \textsuperscript{(54)}, \VUgras{ਬੱਕਰੀਆਂ} \textsuperscript{(51)} ਤੇ
\VUgras{ਬੱਕਰੀ ਦੇ ਬੱਚੇ} \textsuperscript{(52)} ਹੱਕਦਾ ਹੈ। ਉਸ ਦਾ ਵਫ਼ਾਦਾਰ
\VUgras{ਕੁੱਤਾ} \textsuperscript{(55)} ਆਰਗਸ ਸਭ ਤੋਂ ਪਿੱਛੇ ਆਉਂਦਾ ਹੈ ਤੇ ਭੌਂਕ ਕੇ
ਪਿਛੜਿਆਂ ਨੂੰ ਅੱਗੇ ਕਰਦਾ ਹੈ।

%%K t07-c2-11-1 p
11. --- ਇਹ ਸਾਰਾ ਰੌਲਾ-ਰੱਪਾ ਉਸ ਭਲੀ \VUgras{ਦੁੱਧ ਵਾਲੀ ਗਾਂ} \textsuperscript{(56)} ਦੀ
ਸ਼ਾਂਤੀ ਨਹੀਂ ਵਿਗਾੜਦਾ, ਜੋ ਆਪਣੀ \VUgras{ਟੱਲੀ} \textsuperscript{(57)} ਵਜਾਉਂਦੀ ਰਹਿੰਦੀ
ਹੈ, ਜਦੋਂ ਉਸ ਨੂੰ \VUgras{ਗਊਸ਼ਾਲਾ} \textsuperscript{(58)} ਦੇ ਬੂਹੇ ਦੇ ਸਾਹਮਣੇ ਚੋਇਆ
ਜਾਂਦਾ ਹੈ।

%%K t07-c2-12-1 p
12. --- ਉਹ ਜਿਵੇਂ ਸਮਝਦੀ ਹੈ ਕਿ ਉਸ ਨੇ ਆਪਣਾ ਦੁੱਧ ਦੇਣਾ ਹੀ ਹੈ, ਤਾਂ ਜੋ
\VUgras{ਖੇਤ ਦੀ ਮਾਲਕਣ} \textsuperscript{(60)} \VUgras{ਮਧਾਣੀ}
\textsuperscript{(61)} ਵਿਚ \VUgras{ਮੱਖਣ} ਬਣਾ ਸਕੇ, ਜਾਂ ਉਹ ਵਧੀਆ \VUgras{ਪਨੀਰ}
\textsuperscript{(64)} ਜੋ \VUgras{ਟੱਬ} \textsuperscript{(63)} ਉੱਤੇ ਰੱਖੇ ਪਟੜੇ ਉੱਤੇ
ਨਿੱਤਰ ਰਹੇ ਹਨ।

%%K t07-c2-13-1 p
13. --- ਸਾਰਾ \VUgras{ਦੁੱਧ-ਘਰ} \textsuperscript{(59)} \VUgras{ਖੇਤ ਦੇ ਮਜ਼ਦੂਰ}
\textsuperscript{(65)} ਦੀ ਬਦੌਲਤ ਬਹੁਤ ਸਾਫ਼ ਰਹਿੰਦਾ ਹੈ, ਜਿਸ ਨੇ ਹੁਣੇ ਉਸ ਵੱਡੀ ਬਾਲਟੀ
ਵਿਚ \VUgras{ਦੁੱਧ ਦੇ ਡੱਬੇ} \textsuperscript{(62)} ਧੋਤੇ ਹਨ ਜੋ ਉਹ ਚੁੱਕੀ ਹੋਇਆ ਹੈ।

%%K t07-tit-9 sub
\VUsaut{3.0mm}
\VUpk{10.2pt}{40}{ਕੁਕੜਖ਼ਾਨਾ।}
\VUsaut{3.0mm}

%%K t07-c2-14-1 p
14. --- ਆਪਣੀਆਂ ਭਾਰੀਆਂ ਖੜਾਵਾਂ ਵਿਚ ਉਹ ਕੁਝ ਭੱਦੇ ਢੰਗ ਨਾਲ ਤੁਰਦਾ ਹੈ, ਤੇ ਇਸ ਨਾਲ
\VUgras{ਟਰਕੀ} \textsuperscript{(68)}, \VUgras{ਬੱਤਖਾਂ} \textsuperscript{(66)},
\VUgras{ਬੱਤਖਾਂ ਦੇ ਨਰ} \textsuperscript{(67)} ਤੇ \VUgras{ਹੰਸ} \textsuperscript{(69)}
ਉੱਡ ਜਾਂਦੇ ਹਨ, ਜੋ ਆਪਣੀਆਂ \VUgras{ਚੁੰਝਾਂ} \textsuperscript{(70)} ਚੌੜੀਆਂ ਖੋਲ੍ਹ ਕੇ ਬੇਚੈਨ
ਰੌਲਾ ਪਾਉਂਦੇ ਹਨ।

%%K t07-c2-14-2 p
\VUgras{ਕੁੱਕੜ} \textsuperscript{(71)}, ਜੋ ਵੱਧ ਲੜਾਕਾ ਹੈ, ਆਪਣੀ ਲਾਲ \VUgras{ਕਲਗੀ}
\textsuperscript{(72)} ਚੁੱਕਦਾ ਹੈ, ਆਪਣੀ \VUgras{ਪੂਛ} \textsuperscript{(73)} ਦੇ ਖੰਭ
ਝਾੜਦਾ ਹੈ ਤੇ ਇਕ ਗੂੰਜਦੀ « ਬਾਂਗ » ਛੱਡਦਾ ਹੈ।

%%K t07-c2-15-1 p
15. --- ਇਕ \VUgras{ਕੁਕੜੀ-ਮਾਂ} \textsuperscript{(74)} \VUgras{ਗੋਹੇ ਦੇ ਢੇਰ}
\textsuperscript{(81)} ਉੱਤੇ ਚੁਗ ਰਹੀ ਹੈ, ਆਪਣੇ \VUgras{ਚੂਚਿਆਂ}
\textsuperscript{(75)} ਨਾਲ। \VUgras{ਸੂਰ ਦਾ ਬੱਚਾ} \textsuperscript{(78)} ਆਪਣੀ ਮਾਂ ਵੱਲ
ਦੌੜ ਜਾਂਦਾ ਹੈ, ਉਸ ਵੱਡੀ \VUgras{ਸੂਰੀ} \textsuperscript{(77)} ਵੱਲ ਜੋ ਸਾਨੂੰ ਵਾੜੇ ਕੋਲ
ਵਿਖਾਈ ਦਿੰਦੀ ਹੈ।

%%K t07-c2-16-1 p
16. --- ਆਪਣੇ ਡੱਬਿਆਂ ਵਿਚ \VUgras{ਸਹੇ} \textsuperscript{(79)} ਉਹ ਕੋਮਲ \VUgras{ਘਾਹ}
\textsuperscript{(80)} ਚਾਅ ਨਾਲ ਖਾ ਰਹੇ ਹਨ ਜੋ ਮਾਲਕ ਦੀ ਧੀ ਉਨ੍ਹਾਂ ਲਈ ਲਿਆਉਂਦੀ ਹੈ। ਜੈਨੀ
ਨੂੰ ਆਪਣੇ ਸਹੇ ਬਹੁਤ ਪਿਆਰੇ ਹਨ, ਤੇ ਹਰ ਦਿਨ, ਡੱਬੇ ਤੋਂ ਤਾਜ਼ੇ \VUgras{ਆਂਡੇ}
\textsuperscript{(76)} ਲਿਆਉਣ ਪਿੱਛੋਂ, ਉਹ ਆਪਣੇ ਨਿੱਕੇ ਮਿੱਤਰਾਂ ਨੂੰ ਮਿਲਣ ਆਉਂਦੀ ਹੈ।
\VUsaut{6.0mm}
\VUfilet{22mm}
